% ============ Solutions and hints for exercises (Part 2: §4-§6) ============
% 原书 pp.288-289; 原书对 §4 的 Ex.4.1-4.4,4.7-4.9,4.12 与 §5 之外及 §6 的 Ex.6.5-6.7、附录 6.8-6.9 未给解答, 编号按原书跳号.

\section*{\S\ref{sec:4}}

\noindent\textbf{4.5.}\quad Write $V(I_\lambda)$ for the complement of $U_\lambda$, where $I_\lambda$ is an ideal of $A$. Then
$\bigcap V(I_\lambda) = V(\sum I_\lambda) = \emptyset$, so $1 \in \sum I_\lambda$, and therefore a finite sum of $I_\lambda$ also
contains $1$.

\noindent\textbf{4.6.}\quad If $\Spec A = V(I_1) \cup V(I_2)$ with $V(I_1) \cap V(I_2) = \emptyset$ then $I_1 + I_2 = A$ and
$I_1I_2 \subset \nil(A)$. So $1 = e_1 + e_2$ with $e_i \in I_i$ for $i = 1, 2$ and $(e_1e_2)^n = 0$. So $1 =
(e_1 + e_2)^{2n} = e_1^{2n}x_1 + e_2^{2n}x_2$ with $x_i \in A$. So $e = e_1^{2n}x_1$ satisfies $e(1 - e) = 0$.

\noindent\textbf{4.10.}\quad For $p \in \Spec A$, if $V(p) = V(a) \cup V(b)$ then $p \in V(p)$ gives $p \supset a$ or $p \supset b$, and
hence either $V(p) = V(a)$ or $V(p) = V(b)$. Conversely, if $V(I)$ is irreducible,
then for $x, y \in A$ with $xy \in \sqrt{I}$, from $V = V(I + Ax) \cup V(I + Ay)$ we have,
say, $V = V(I + Ax)$, and $x \in \sqrt{I}$, this proves $\sqrt{I} \in \Spec A$.

\noindent\textbf{4.11.}\quad If there is a closed subset which cannot be so expressed, let $V$ be a minimal
one. Then $V$ must be reducible, but if $V = V_1 \cup V_2$ with $V_i \neq V$ then, by
minimality, each of $V_1$ and $V_2$ is a union of a finite number of irreducible
closed sets, hence also $V$, a contradiction.

\section*{\S\ref{sec:5}}

\noindent\textbf{5.1.}\quad Set $k[X_1, \ldots, X_n]/p = k[x_1, \ldots, x_n]$; then by Theorem~\hyperref[thm:5.6]{6}, $\coht p =
\trdeg_k k(x)$. Suppose this is $r$, and that $x_1, \ldots, x_r$ is a transcendence basis
of $k(x)$ over $k$, and set $K = k(X_1, \ldots, X_r)$; then $k[X_1, \ldots, X_n]_p$ is the
localisation of $K[X_{r+1}, \ldots, X_n]$ at a prime ideal $P$, with $\htop p = \htop P$. This
reduces us to proving that if $r = 0$ then $\htop p = n$. In this case $x_1, \ldots, x_n$ are
algebraic over $k$, and letting $p_i$ be the kernel of $k[X_1, \ldots, X_n] \to
k[x_1, \ldots, x_i, X_{i+1}, \ldots, X_n]$ we get a strictly increasing chain $0 \subset p_1 \subset p_2
\subset \cdots \subset p_n = p$, giving $\htop p \ge n$, but by the corollary to Theorem~\hyperref[thm:5.6]{6},
$\htop p \le n$.

\noindent\textbf{5.2.}\quad If $A$ is a zero-dimensional Noetherian ring then all prime ideals are
minimal, and by Ex.~\ref{ex:4.12-ex}, there are only finitely many of these. Let these be
$p_1, \ldots, p_n$; then since $p_1 \cdots p_n = \nil(A)$ there is an $n$ such that $(p_1 \cdots p_n)^n = 0$.
For any ideal $I$ and any $i$, the module $I/Ip_i$ is a finite-dimensional vector
space over $A/p_i$, so that $l(I/Ip_i) < \infty$. It follows easily that $l(A) < \infty$, so that
$A$ is Artinian.

\section*{\S\ref{sec:6}}

\noindent\textbf{6.1.}\quad $\Ass M = \{(0), (3)\}$. ($\supset$ is obvious, $\subset$ from Theorem~\hyperref[thm:6.3]{3}.)

\noindent\textbf{6.2.}\quad No. Let $M$ be as in the previous question, $M_1 = \{(a, 0) \mid a \in \Z\}$ and
$M_2 = \{(0, a) \mid a \in \Z\}$; then $M = M_1 + M_2$, but each $M_i \cong \Z$.

\noindent\textbf{6.3.}\quad Since $xA/x^nA \cong A/x^{n-1}A$, there is an exact sequence
\[
  0 \to A/x^{n-1}A \to A/x^nA \to A/xA \to 0.
\]

\noindent\textbf{6.4.}\quad Use a primary decomposition of $I$.
